\chapter{Green Stool}

\section*{Definition}
Green-colored stool resulting from the presence of bile pigments, rapid intestinal transit, or ingestion of green-colored foods or supplements.

\section*{Pathophysiology}
Bilirubin produced from heme breakdown is conjugated in the liver, secreted into bile, and normally metabolized by gut bacteria to stercobilin (brown pigment). Green stool occurs when bile passes through the intestine too rapidly for bacterial conversion to stercobilin, or when biliverdin (green pigment) is not fully reduced. Dietary chlorophyll (green vegetables) and artificial food coloring (green dyes) are common benign causes.

\section*{Etiology}
\begin{itemize}
  \item Dietary: green leafy vegetables (spinach, kale)
  \item Food coloring (green dyes in beverages, candies, gelatin)
  \item Rapid intestinal transit (diarrhea, IBS-D)
  \item Bile acid malabsorption
  \item Antibiotic use (altered gut flora)
  \item Iron supplements
  \item Green-colored medications or supplements
  \item Infectious gastroenteritis (especially Salmonella)
  \item Crohn disease or ulcerative colitis
  \item Malabsorption syndromes
\end{itemize}

\section*{Indications for Evaluation}
Seek care if:
\begin{itemize}
  \item Severe or worsening symptoms
  \item Persistent green stool with diarrhea, weight loss, abdominal pain, or other concerning GI symptoms
  \item Accompanied by other concerning signs
\end{itemize}

\section*{Related Symptoms}
\begin{itemize}
  \item Diarrhea
  \item Abdominal pain
  \item Bloating
  \item Weight loss
  \item Nausea
\end{itemize}

\textbf{Note:} Always consider serious underlying conditions when evaluating persistent green stool.

\section*{Self-Care / Home Management}
\begin{itemize}
  \item Rest and avoid activities that may aggravate the condition.
  \item Maintain adequate hydration and a balanced diet.
  \item Over-the-counter remedies may provide symptomatic relief (consult a pharmacist or doctor).
  \item Monitor symptoms and seek professional advice if they persist or worsen.
  \item Avoid self-medicating with prescription medications without medical guidance.
\end{itemize}

\section*{Diagnostic Approach}
\begin{itemize}
  \item A thorough history and physical examination are the first steps in evaluation.
  \item Laboratory tests (blood work, urinalysis) may be ordered based on clinical suspicion.
  \item Imaging studies (X-ray, ultrasound, CT, MRI) may be indicated when structural pathology is suspected.
  \item Specialist referral is arranged when the diagnosis remains uncertain or requires specific expertise.
\end{itemize}

\section*{Treatment / Management Overview}
\begin{itemize}
  \item Treatment targets the underlying cause identified through diagnostic evaluation.
  \item Supportive care and symptom management are provided while awaiting definitive therapy.
  \item Medications, lifestyle modifications, and physical therapies may be prescribed as appropriate.
  \item Follow-up ensures treatment effectiveness and allows timely adjustments to the care plan.
\end{itemize}

\clearpage